\documentclass[11pt,a4paper]{article}
\usepackage[utf8]{inputenc}
\usepackage[T1]{fontenc}
\usepackage[ngerman]{babel}
\usepackage{helvet}
\renewcommand{\familydefault}{\sfdefault}
\usepackage{geometry}
\usepackage{hyperref}
\usepackage{booktabs}
\usepackage{longtable}
\usepackage{array}
\usepackage{enumitem}
\usepackage{xcolor}
\usepackage{fancyhdr}
\usepackage{titlesec}
\usepackage{tcolorbox}
\usepackage{tikz}
\usepackage{colortbl}
\usepackage{tabularx}
\usepackage{multirow}
\usepackage{graphicx}
\usepackage{amssymb}
\usepackage{pifont}
\usepackage{pdflscape}

\usetikzlibrary{shapes.geometric, arrows.meta, positioning, calc, fit, backgrounds, shapes.symbols}

\geometry{margin=1.8cm, top=2.2cm, bottom=1.8cm}
\hypersetup{colorlinks=true,linkcolor=mainblue,urlcolor=mainblue}

% Farben
\definecolor{mainblue}{RGB}{23,74,108}
\definecolor{lightblue}{RGB}{220,235,245}
\definecolor{accentgreen}{RGB}{34,139,94}
\definecolor{accentorange}{RGB}{230,126,34}
\definecolor{lightgray}{RGB}{245,245,245}
\definecolor{rp2002}{RGB}{180,180,180}
\definecolor{rp2025}{RGB}{34,139,94}
\definecolor{thema1}{RGB}{70,130,180}
\definecolor{thema2}{RGB}{139,90,43}
\definecolor{thema3}{RGB}{178,34,34}
\definecolor{thema4}{RGB}{60,120,60}
\definecolor{examcolor}{RGB}{128,0,128}

% Symbole
\newcommand{\cmark}{\textcolor{accentgreen}{\ding{51}}}
\newcommand{\xmark}{\textcolor{red}{\ding{55}}}
\newcommand{\lmark}{\textcolor{accentorange}{\textbf{[LK]}}}
\newcommand{\checkbox}{$\square$}

% Section Styling
\titleformat{\section}{\Large\bfseries\color{mainblue}}{\thesection.}{0.5em}{}[\vspace{-0.5em}\textcolor{mainblue}{\titlerule[1pt]}]
\titleformat{\subsection}{\large\bfseries\color{mainblue}}{\thesubsection}{0.5em}{}
\titleformat{\subsubsection}{\normalsize\bfseries\color{mainblue}}{\thesubsubsection}{0.5em}{}

% Header/Footer
\pagestyle{fancy}
\fancyhf{}
\fancyhead[L]{\small\color{gray}Leitfaden Spanisch MV v4}
\fancyhead[R]{\small\color{gray}\thepage}
\renewcommand{\headrulewidth}{0.4pt}

% Boxen
\newtcolorbox{infobox}[1][]{
  colback=lightblue, colframe=mainblue, fonttitle=\bfseries, title=#1,
  boxrule=0.5pt, arc=2pt, left=5pt, right=5pt, top=5pt, bottom=5pt
}
\newtcolorbox{wichtigbox}{
  colback=lightblue, colframe=mainblue, boxrule=0pt, leftrule=3pt, arc=0pt,
  left=8pt, right=5pt, top=5pt, bottom=5pt
}
\newtcolorbox{themenbox}[2][]{
  colback=#2!10, colframe=#2, fonttitle=\bfseries, title=#1,
  boxrule=1pt, arc=2pt, left=5pt, right=5pt, top=5pt, bottom=5pt
}
\newtcolorbox{literaturbox}{
  colback=lightgray, colframe=gray, boxrule=0.5pt, arc=2pt,
  left=5pt, right=5pt, top=3pt, bottom=3pt
}
\newtcolorbox{checklistbox}[1][]{
  colback=accentgreen!10, colframe=accentgreen, fonttitle=\bfseries, title=#1,
  boxrule=1pt, arc=2pt, left=5pt, right=5pt, top=5pt, bottom=5pt
}
\newtcolorbox{exambox}[1][]{
  colback=examcolor!10, colframe=examcolor, fonttitle=\bfseries, title=#1,
  boxrule=1pt, arc=2pt, left=5pt, right=5pt, top=5pt, bottom=5pt
}
\newtcolorbox{fallbeispielbox}[1][]{
  colback=accentorange!10, colframe=accentorange, fonttitle=\bfseries, title=#1,
  boxrule=1pt, arc=2pt, left=5pt, right=5pt, top=5pt, bottom=5pt
}

\newcolumntype{H}{>{\columncolor{lightblue}\bfseries}l}
\newcolumntype{C}{>{\columncolor{lightblue}\bfseries}c}

\begin{document}

% === TITELSEITE ===
\begin{titlepage}
\centering
\vspace*{2.5cm}
{\Huge\bfseries\color{mainblue} Spanisch in\\[0.3em]Mecklenburg-Vorpommern\par}
\vspace{0.6cm}
{\LARGE\bfseries Umfassender Leitfaden v4\par}
\vspace{0.4cm}
{\large\textit{Mit Checklisten, Prüfungsvorbereitung und Quick-Reference}\par}
\vspace{0.5cm}
{\large Rahmenpläne \textbullet{} Prüfungen \textbullet{} Themenfelder \textbullet{} Literatur\par}
\vspace{1.5cm}

\begin{tcolorbox}[colback=lightgray, colframe=lightgray, width=13cm, boxrule=0pt]
\textbf{Stand:} Dezember 2025\\[0.2em]
\textbf{Grundlage:} 22 offizielle Dokumente des Ministeriums für Bildung M-V\\[0.2em]
\textbf{Geltungsbereich:} Gymnasium, Regionale Schule, Gymnasiale Oberstufe\\[0.2em]
\textbf{NEU in v4:} Checklisten, Bewertungsbögen, Examenshinweise, Quick-Reference
\end{tcolorbox}

\vspace{1.5cm}
\begin{tikzpicture}
\node[rectangle, rounded corners=5pt, fill=accentgreen!20, draw=accentgreen, line width=1pt, minimum width=10cm, minimum height=1.2cm] {
  \textbf{Zielgruppe:} Lehrkräfte \textbullet{} Referendar:innen \textbullet{} Lehramtsstudierende
};
\end{tikzpicture}

\vfill
{\small Erstellt auf Grundlage der offiziellen Dokumente des IQ M-V\par}
\end{titlepage}

% === INHALTSVERZEICHNIS ===
\tableofcontents
\newpage

% ============================================================================
% SCHNELLZUGRIFF: ENTSCHEIDUNGS-FLOWCHART
% ============================================================================
\section{Schnellzugriff: Welcher Rahmenplan gilt?}

\begin{center}
\begin{tikzpicture}[
  decision/.style={diamond, draw=mainblue, fill=lightblue, text width=2.8cm, text centered, aspect=2, font=\small},
  block/.style={rectangle, rounded corners=3pt, draw=mainblue, fill=lightblue, text width=3.5cm, text centered, minimum height=1cm, font=\small},
  result/.style={rectangle, rounded corners=3pt, draw=accentgreen, fill=accentgreen!20, text width=4cm, text centered, minimum height=1cm, font=\small\bfseries},
  arrow/.style={->, thick, >=stealth}
]

% Start
\node[block] (start) at (0,0) {Spanischunterricht\\planen};

% Entscheidung 1: Sek I oder II?
\node[decision, below=0.8cm of start] (sek) {Sek I oder\\Sek II?};

% Sek II Pfad
\node[result, right=2.5cm of sek] (sek2) {RP Sek II (2019)\\+ Operatorenliste};

% Sek I: Welches Schuljahr?
\node[decision, below=1.2cm of sek] (jahr) {Schuljahr\\2025/26+?};

% Vor 2025
\node[block, left=2cm of jahr] (alt) {RP 2002\\(auslaufend)};

% Ab 2025: Welche Jg?
\node[decision, below=1.2cm of jahr] (jg) {Jg. 7 neu\\gestartet?};

% Ergebnisse
\node[result, left=2cm of jg] (rp2002) {RP 2002\\für diese Klasse};
\node[result, right=2cm of jg] (rp2025) {RP 2025\\(Gym oder RegS)};

% Schulart
\node[decision, below=1cm of rp2025] (schulart) {Schulart?};
\node[result, below left=0.8cm and 0.5cm of schulart] (gym) {RP Gym 2025\\Ziel: B1};
\node[result, below right=0.8cm and 0.5cm of schulart] (regs) {RP RegS 2025\\Ziel: A2+};

% Pfeile
\draw[arrow] (start) -- (sek);
\draw[arrow] (sek) -- node[above] {Sek II} (sek2);
\draw[arrow] (sek) -- node[left] {Sek I} (jahr);
\draw[arrow] (jahr) -- node[above] {Nein} (alt);
\draw[arrow] (jahr) -- node[left] {Ja} (jg);
\draw[arrow] (jg) -- node[above] {Nein} (rp2002);
\draw[arrow] (jg) -- node[above] {Ja} (rp2025);
\draw[arrow] (rp2025) -- (schulart);
\draw[arrow] (schulart) -- node[left, font=\scriptsize] {Gym/IGS} (gym);
\draw[arrow] (schulart) -- node[right, font=\scriptsize] {RegS} (regs);
\end{tikzpicture}
\end{center}

\begin{wichtigbox}
\textbf{Merke:} Ab 2025/26 gilt der neue RP \textbf{nur für Jg. 7}. Höhere Jahrgänge arbeiten weiter mit RP 2002, bis sie die Schule verlassen. Erst ab \textbf{2028/29} gilt RP 2025 für alle!
\end{wichtigbox}

% ============================================================================
% TEIL I: GRUNDLAGEN
% ============================================================================
\newpage
\part*{Teil I: Grundlagen}
\addcontentsline{toc}{part}{Teil I: Grundlagen}

\section{Dokumentenübersicht}

\subsection{Rahmenpläne}

\begin{tabularx}{\textwidth}{H X c c c}
\rowcolor{lightblue}
\textbf{Dokument} & \textbf{Schulart} & \textbf{Jg.} & \textbf{Ab} & \textbf{Status} \\
\hline
RP Spanisch Gym 2025 & Gymnasium, IGS & 7--10 & 2025/26 & \textbf{NEU} \\
\rowcolor{lightgray}
RP Spanisch RegS 2025 & Regionale Schule & 7--10 & 2025/26 & \textbf{NEU} \\
RP Spanisch Sek II & Gym. Oberstufe & 10--12 & 2019/20 & GÜLTIG \\
\rowcolor{lightgray}
RP Spanisch Gym 2002 & Gym. + IGS & 7--10 & 2002 & AUSLAUFEND \\
\hline
\end{tabularx}

\subsection{Begleitdokumente}

\begin{tabularx}{\textwidth}{H X c}
\rowcolor{lightblue}
\textbf{Dokument} & \textbf{Inhalt} & \textbf{S.} \\
\hline
Operatorenliste & AFB I--III Definitionen (verbindlich für Abitur!) & 6 \\
\rowcolor{lightgray}
Stilmittel-Liste & Rhetorische Mittel für Textanalyse & 1 \\
Literaturempfehlungen & Kommentierte Lektüren für Oberstufe & 3 \\
\rowcolor{lightgray}
Verknüpfungsbeispiele (4×) & Unterrichtsbeispiele zu allen 4 Themen & 8 \\
El método Grönholm & Annotiertes Theaterstück (LK-Empfehlung) & 99 \\
\rowcolor{lightgray}
HR Sprechprüfung & Komplexe Leistungsermittlung Sprechen & 13 \\
HR Paarprüfung & Mündl. Abiturprüfung als Paarformat (ab 2025) & 9 \\
\rowcolor{lightgray}
RP Sprachbildung & Fachübergreifende Sprachförderung & 89 \\
\hline
\end{tabularx}

\section{Gültigkeit und Übergänge}

\begin{center}
\begin{tikzpicture}[scale=0.8]
  \fill[rp2002] (0,5.3) rectangle (0.5,5.6); \node[right] at (0.6,5.45) {\small RP 2002};
  \fill[rp2025] (3.5,5.3) rectangle (4,5.6); \node[right] at (4.1,5.45) {\small RP 2025};

  \node[left] at (-0.3,4) {\small\textbf{Jg. 7:}};
  \node[left] at (-0.3,3) {\small\textbf{Jg. 8:}};
  \node[left] at (-0.3,2) {\small\textbf{Jg. 9:}};
  \node[left] at (-0.3,1) {\small\textbf{Jg. 10:}};

  \node at (1,0) {\small 24/25}; \node at (3.5,0) {\small 25/26};
  \node at (6,0) {\small 26/27}; \node at (8.5,0) {\small 27/28}; \node at (11,0) {\small 28/29};

  \draw[gray,dashed] (2.25,0.5) -- (2.25,4.5);
  \draw[gray,dashed] (4.75,0.5) -- (4.75,4.5);
  \draw[gray,dashed] (7.25,0.5) -- (7.25,4.5);
  \draw[gray,dashed] (9.75,0.5) -- (9.75,4.5);

  \fill[rp2002] (0,3.7) rectangle (2.25,4.3); \fill[rp2025] (2.25,3.7) rectangle (12,4.3);
  \fill[rp2002] (0,2.7) rectangle (4.75,3.3); \fill[rp2025] (4.75,2.7) rectangle (12,3.3);
  \fill[rp2002] (0,1.7) rectangle (7.25,2.3); \fill[rp2025] (7.25,1.7) rectangle (12,2.3);
  \fill[rp2002] (0,0.7) rectangle (9.75,1.3); \fill[rp2025] (9.75,0.7) rectangle (12,1.3);
\end{tikzpicture}
\end{center}

\subsection{Niveaustufen im Vergleich}

\begin{center}
\begin{tabular}{l|c|c|c|c}
\rowcolor{lightblue}
\textbf{Jg.} & \textbf{Gym USt} & \textbf{Gym Niveau} & \textbf{RegS USt} & \textbf{RegS Niveau} \\
\hline
7 & 120 & A1 & 90 & A1 \\
\rowcolor{lightgray}
8 & 120 & A2 & 90 & A1+ \\
9 & 90 & A2+ & 90 & A2 \\
\rowcolor{lightgray}
10 & 90 & \textbf{B1} & 90 & \textbf{A2+} \\
\hline
\multicolumn{5}{l}{\footnotesize Sek II: Einführungsphase B1/A2, Qualifikationsphase \textbf{B2}/B1(+)} \\
\end{tabular}
\end{center}

% ============================================================================
% TEIL II: VERBINDLICHE INHALTE SEK I (DETAILLIERT)
% ============================================================================
\newpage
\section{Verbindliche Inhalte Sekundarstufe I -- Detailliert}

\begin{center}
\begin{tabularx}{\textwidth}{c|l|X}
\rowcolor{lightblue}
\textbf{Jg.} & \textbf{Thema} & \textbf{Inhalte und Schwerpunkte} \\
\hline
\multirow{3}{*}{7} & yo y mi entorno & Familie, Freunde, Haustiere, Schule, Klassenzimmer \\
& ¿Dónde y cómo vivo? & Wohnung, Zimmer, Stadtviertel, Wegbeschreibung \\
& mi día & Tagesablauf, Uhrzeiten, Wochentage, Schulalltag \\
\hline
\rowcolor{lightgray}
\multirow{3}{*}{8} & mi tiempo libre & Hobbys, Sport, Musik, Verabredungen \\
\rowcolor{lightgray}
& vivir en armonía & Umwelt, Nachhaltigkeit, Zusammenleben \\
\rowcolor{lightgray}
& mundo hispanohablante & Länder, Feste, Traditionen, Geografie \\
\hline
\multirow{3}{*}{9} & la vida cotidiana & Einkaufen, Essen, Restaurant, Mode \\
& Barcelona & Sehenswürdigkeiten, Geschichte, Architektur (Gaudí) \\
& un/a artista & Künstlerbiografie (Picasso, Dalí, Kahlo, etc.) \\
\hline
\rowcolor{lightgray}
\multirow{3}{*}{10} & España -- contrastes & Regionen, Autonomien, Nord-Süd-Unterschiede \\
\rowcolor{lightgray}
& el mundo laboral & Berufe, Bewerbung, Praktikum, Zukunftspläne \\
\rowcolor{lightgray}
& retos del siglo XXI & Migration, Digitalisierung, Klimawandel \\
\hline
\end{tabularx}
\end{center}

\begin{wichtigbox}
\textbf{80/20-Regel:} 80\% verbindliche Inhalte, 20\% Freiräume. \textbf{Kein schulinterner Lehrplan} mehr nötig (RP 2025)!
\end{wichtigbox}

% ============================================================================
% TEIL III: THEMENFELDER OBERSTUFE
% ============================================================================
\newpage
\part*{Teil II: Themenfelder Oberstufe}
\addcontentsline{toc}{part}{Teil II: Themenfelder Oberstufe}

\section{Übersicht der 4 Themenfelder}

\begin{center}
\begin{tikzpicture}[
  thema/.style={rectangle, rounded corners=3pt, minimum width=5.5cm, minimum height=1cm, text centered, font=\small},
]
\node[thema, fill=thema1!25, draw=thema1, line width=1pt] (t1) at (0,0) {\textbf{1. El individuo y la sociedad}\\45 USt (GK) / 75 USt (LK)};
\node[thema, fill=thema2!25, draw=thema2, line width=1pt, right=0.4cm of t1] (t2) {\textbf{2. La identidad nacional}\\45 USt (GK) / 75 USt (LK)};
\node[thema, fill=thema3!25, draw=thema3, line width=1pt, below=0.4cm of t1] (t3) {\textbf{3. Aspectos de política}\\30 USt (GK) / 50 USt (LK)};
\node[thema, fill=thema4!25, draw=thema4, line width=1pt, right=0.4cm of t3] (t4) {\textbf{4. Desafíos modernos}\\30 USt (GK) / 50 USt (LK)};
\end{tikzpicture}
\end{center}

% --- THEMA 1 ---
\subsection{Thema 1: El individuo y la sociedad}

\begin{themenbox}[Schwerpunkte]{thema1}
La situación de la juventud \textbullet{} Facetas de la vida \textbullet{} Formas de convivencia
\end{themenbox}

\textbf{Kompetenzverknüpfung (Kurzfilm „Física II"):}
\begin{itemize}[leftmargin=1cm, itemsep=1pt]
\item \textbf{Hörsehverstehen:} Hauptaussagen entnehmen; Jorges Beweggründe analysieren
\item \textbf{Lesen:} Interview mit Regisseur; Intentionen herausarbeiten
\item \textbf{Sprechen:} Diskussion zu Zukunftsperspektiven nach der Schulzeit
\item \textbf{Schreiben:} Tagebucheintrag aus Jorges Sicht
\item \lmark{} \textbf{Zusätzlich:} Rezension schreiben; Atmosphäre-Analyse
\end{itemize}

\begin{literaturbox}
\textbf{Literatur:} Desconexión total (Miquel) \textbullet{} La Casa (Roca, Novela Gráfica) \textbullet{} La mirada (Madrid)\\
\lmark{} \textbf{El método Grönholm} (Galceran) -- siehe Detailübersicht S. \pageref{sec:gronholm}
\end{literaturbox}

% --- THEMA 2 ---
\subsection{Thema 2: La identidad nacional y la diversidad cultural}

\begin{themenbox}[Schwerpunkte]{thema2}
Momentos cruciales \textbullet{} El multiculturalismo \textbullet{} España después de la Transición
\end{themenbox}

\textbf{Kompetenzverknüpfung (Fidel Castro Rede UNO 1979):}
\begin{itemize}[leftmargin=1cm, itemsep=1pt]
\item \textbf{Hörsehverstehen:} Politische Rede verstehen; Hauptaussagen entnehmen
\item \textbf{Lesen:} Redetext lesen; Absicht und Wirkung erfassen
\item \textbf{Sprechen:} Vortrag zur politischen Entwicklung eines lat.am. Landes
\item \textbf{Schreiben:} Nachricht/Bericht zu gesellschaftspolitischem Aspekt
\item \lmark{} \textbf{Zusätzlich:} Eigene Rede verfassen; implizite Informationen
\end{itemize}

\begin{literaturbox}
\textbf{Literatur:} El desertador (Merino) \textbullet{} Una fosa poco común (Jambrina) \textbullet{} Las venas abiertas de América Latina (Galeano)
\end{literaturbox}

% --- THEMA 3 ---
\subsection{Thema 3: Aspectos actuales de la política y de la sociedad}

\begin{themenbox}[Schwerpunkte]{thema3}
La globalización \textbullet{} Responsabilidades sociales \textbullet{} Movimientos migratorios
\end{themenbox}

\textbf{Kompetenzverknüpfung (La casa en Mango Street):}
\begin{itemize}[leftmargin=1cm, itemsep=1pt]
\item \textbf{Hörsehverstehen:} Doku „México--EE.UU.: El gran cruce"
\item \textbf{Lesen:} Kapitel aus „La casa en Mango Street" (Sandra Cisneros)
\item \textbf{Sprechen:} Rollenspiel: Lateinamerikanische Emigranten in USA
\item \textbf{Schreiben:} Brief aus Sicht Esperanzas
\item \lmark{} \textbf{Zusätzlich:} Zusätzliches Kapitel verfassen; Varietäten des Español Americano
\end{itemize}

\begin{literaturbox}
\textbf{Literatur:} La piel de un indio no cuesta caro (Ribeyro) \textbullet{} Abdel (Páez)
\end{literaturbox}

% --- THEMA 4 ---
\subsection{Thema 4: Desafíos modernos}

\begin{themenbox}[Schwerpunkte]{thema4}
El mundo en movimiento \textbullet{} Sueños y realidades \textbullet{} Tendencias ecológicas
\end{themenbox}

\textbf{Kompetenzverknüpfung (Rigoberta Menchú):}
\begin{itemize}[leftmargin=1cm, itemsep=1pt]
\item \textbf{Hörsehverstehen:} Pressekonferenz zu „El agua y la humanidad"
\item \textbf{Lesen:} Auszüge aus „El agua y la humanidad"
\item \textbf{Sprechen:} Diskussion: Maßnahmen zur Wasserversorgung
\item \textbf{Schreiben:} Artikel für Schülermagazin
\item \lmark{} \textbf{Zusätzlich:} Offener Brief; zielorientiert argumentieren
\end{itemize}

\begin{literaturbox}
\textbf{Literatur:} 547 amigos (Silva) \textbullet{} \lmark{} El guardagujas (Arreola)
\end{literaturbox}

% ============================================================================
% EL MÉTODO GRÖNHOLM -- DETAILÜBERSICHT
% ============================================================================
\newpage
\section{El método Grönholm -- Detailübersicht}
\label{sec:gronholm}

\begin{infobox}[Jordi Galceran, 2003 -- Theaterstück für Leistungskurs]
\textbf{Thema:} Assessment-Center für einen hochdotierten Posten in einer Schweizer Firma.\\
\textbf{Themenanbindung:} El individuo y la sociedad (Arbeitswelt, ethische Fragen)\\
\textbf{Quelle:} Kostenlos annotiert auf \url{https://www.bildung-mv.de/.../spanisch/}
\end{infobox}

\subsection{Szenenübersicht}

\begin{center}
\begin{tabularx}{\textwidth}{c|l|X}
\rowcolor{lightblue}
\textbf{Szene} & \textbf{Titel/Fokus} & \textbf{Inhalt} \\
\hline
1 & Llegada & 4 Bewerber:innen treffen ein; erste Spannungen \\
\rowcolor{lightgray}
2 & Primera prueba & Aufgabe: Wer ist der/die Maulwurf/in der Firma? \\
3 & Descubrimiento & Carlos wird als Komplize entlarvt und ausgeschlossen \\
\rowcolor{lightgray}
4 & Segunda prueba & Persönliche Geheimnisse müssen offenbart werden \\
5 & Conflicto & Moralische Grenzen werden überschritten \\
\rowcolor{lightgray}
6 & Tercera prueba & Finale Entscheidung; psychologischer Druck steigt \\
7 & Desenlace & Überraschende Wendung; Frage: War es das wert? \\
\hline
\end{tabularx}
\end{center}

\subsection{Unterrichtsvorschläge}

\begin{itemize}[leftmargin=1cm]
\item \textbf{Rollenspiel:} SuS übernehmen Bewerberrollen; Szenen nachspielen
\item \textbf{Diskussion:} Ethische Grenzen bei Bewerbungsverfahren
\item \textbf{Schreiben:} Alternatives Ende; Bewerbungsbrief aus Sicht einer Figur
\item \textbf{Analyse:} Sprachliche Mittel der Manipulation
\end{itemize}

% ============================================================================
% TEIL III: PRÜFUNGEN
% ============================================================================
\newpage
\part*{Teil III: Prüfungsformate}
\addcontentsline{toc}{part}{Teil III: Prüfungsformate}

\section{Übersicht}

\begin{center}
\begin{tabularx}{\textwidth}{X|c|c|c}
\rowcolor{lightblue}
\textbf{Format} & \textbf{Rechtsgrundlage} & \textbf{Ab} & \textbf{Pflicht?} \\
\hline
Komplexe Leistungsermittlung Sprechen & § 22 Abs. 5 APVO & 2019 & \textbf{JA} \\
\rowcolor{lightgray}
Klausurersatz durch Sprechprüfung & § 22 Abs. 4 APVO & 2019 & Nein (max. 1/3) \\
Paarprüfung mündl. Prüfungsteil & § 25 Abs. 8 APVO & 2025 & Nein \\
\rowcolor{lightgray}
Paarprüfung mündl. Abitur & § 38 Abs. 6 APVO & 2025 & Nein (Antrag) \\
\hline
\end{tabularx}
\end{center}

\section{Paarprüfung: Ablauf und Bewertung}

\begin{center}
\begin{tikzpicture}[
  block/.style={rectangle, rounded corners=2pt, minimum width=3.2cm, minimum height=0.9cm, text centered, font=\small},
  phase/.style={block, fill=lightblue, draw=mainblue},
  info/.style={block, fill=lightgray, draw=gray, minimum width=5cm}
]
\node[phase] (t1) at (0,0) {\textbf{Teil 1}\\Aufwärmphase};
\node[right=0.1cm of t1, font=\footnotesize] {2'};
\node[info, right=1.2cm of t1] {Interview (KEINE Bewertung)};

\node[phase, below=0.3cm of t1] (t2) {\textbf{Teil 2}\\Monolog};
\node[right=0.1cm of t2, font=\footnotesize] {4'/P};
\node[info, right=1.2cm of t2] {Impulsgesteuertes Sprechen};

\node[phase, below=0.3cm of t2] (t3) {\textbf{Teil 3}\\Dialog};
\node[right=0.1cm of t3, font=\footnotesize] {5'/P};
\node[info, right=1.2cm of t3] {Diskussion (anderes Thema!)};

\node[block, fill=accentgreen!20, draw=accentgreen, below=0.3cm of t3] (bew) {\textbf{Bewertung}};
\node[right=0.1cm of bew, font=\footnotesize] {10'};
\node[info, right=1.2cm of bew] {\textbf{60\% Sprache / 40\% Inhalt}};

\draw[->, thick, gray] (t1) -- (t2);
\draw[->, thick, gray] (t2) -- (t3);
\draw[->, thick, gray] (t3) -- (bew);
\end{tikzpicture}
\end{center}

\subsection{Bewertungsbogen Sprechprüfung (Vorlage)}

\begin{center}
\small
\begin{tabularx}{\textwidth}{|l|X|c|c|c|c|c|}
\hline
\rowcolor{lightblue}
\textbf{Kriterium} & \textbf{Beschreibung} & \textbf{5} & \textbf{4} & \textbf{3} & \textbf{2} & \textbf{1} \\
\hline
\multicolumn{7}{|l|}{\textbf{SPRACHE (60\%)}} \\
\hline
Aussprache & Korrekt, verständlich, Intonation & \checkbox & \checkbox & \checkbox & \checkbox & \checkbox \\
\hline
Wortschatz & Umfang, Präzision, Themenpassung & \checkbox & \checkbox & \checkbox & \checkbox & \checkbox \\
\hline
Grammatik & Korrektheit, Komplexität & \checkbox & \checkbox & \checkbox & \checkbox & \checkbox \\
\hline
Flüssigkeit & Sprechtempo, Pausen, Hesitation & \checkbox & \checkbox & \checkbox & \checkbox & \checkbox \\
\hline
\multicolumn{7}{|l|}{\textbf{INHALT (40\%)}} \\
\hline
Aufgabenerfüllung & Vollständigkeit, Relevanz & \checkbox & \checkbox & \checkbox & \checkbox & \checkbox \\
\hline
Argumentation & Logik, Struktur, Begründung & \checkbox & \checkbox & \checkbox & \checkbox & \checkbox \\
\hline
Interaktion & Eingehen auf Partner:in, Reaktion & \checkbox & \checkbox & \checkbox & \checkbox & \checkbox \\
\hline
\end{tabularx}
\end{center}

\section{Zeitplanung eintägige Prüfung}

\begin{center}
\begin{tabular}{c|c|l}
\rowcolor{lightblue}
\textbf{Zeit} & \textbf{Tandem} & \textbf{Hinweis} \\
\hline
08:00--08:30 & 1 & \\
08:30--09:00 & 2 & \\
09:00--09:30 & 3 & \\
09:30--10:00 & 4 & \\
\rowcolor{accentorange!20}
10:00--10:15 & -- & \textit{Pause 15 Min.} \\
10:15--10:45 & 5 & \\
... & ... & \\
\rowcolor{accentorange!20}
12:30--13:00 & -- & \textit{Mittagspause 30 Min.} \\
... & ... & \\
\textbf{14:15} & \textbf{11} & \textbf{Ende (22 SuS)} \\
\hline
\end{tabular}
\end{center}

\textbf{Bei 3er-Gruppen:} 40 Min./Gruppe (30' Prüfung + 10' Bewertung)

% ============================================================================
% TEIL IV: OPERATOREN
% ============================================================================
\newpage
\part*{Teil IV: Operatoren und AFB}
\addcontentsline{toc}{part}{Teil IV: Operatoren und AFB}

\section{Anforderungsbereiche}

\begin{center}
\begin{tikzpicture}[
  block/.style={rectangle, rounded corners=3pt, minimum width=3.5cm, minimum height=1.2cm, text centered},
]
\node[block, fill=accentgreen!20, draw=accentgreen, line width=1pt] (afb1) at (0,0) {\textbf{AFB I}\\Reproduktion\\25--30\%};
\node[block, fill=accentorange!20, draw=accentorange, line width=1pt, right=0.3cm of afb1] (afb2) {\textbf{AFB II}\\Transfer\\40--50\%};
\node[block, fill=lightblue, draw=mainblue, line width=1pt, right=0.3cm of afb2] (afb3) {\textbf{AFB III}\\Reflexion\\20--30\%};

\node[below=0.15cm of afb1, font=\footnotesize, text width=3.5cm, align=center] {describir, nombrar,\\presentar, resumir};
\node[below=0.15cm of afb2, font=\footnotesize, text width=3.5cm, align=center] {analizar, caracterizar,\\comparar, explicar};
\node[below=0.15cm of afb3, font=\footnotesize, text width=3.5cm, align=center] {comentar, discutir,\\evaluar, justificar};
\end{tikzpicture}
\end{center}

\section{Verbindliche Operatoren}

\begin{center}
\footnotesize
\begin{tabularx}{\textwidth}{c|l|X}
\rowcolor{lightblue}
\textbf{AFB} & \textbf{Operator} & \textbf{Definition} \\
\hline
I & describir & Sachverhalte strukturiert und sachlich darstellen \\
\rowcolor{lightgray}
I & resumir & Textinhalt auf wesentliche Aussagen konzentrieren \\
I & presentar & Sachverhalte/Positionen strukturiert darlegen \\
\rowcolor{lightgray}
I & nombrar & Informationen ohne Erläuterung aufzählen \\
\hline
II & analizar & Materialien systematisch untersuchen und auswerten \\
\rowcolor{lightgray}
II & caracterizar & Eigenarten herausarbeiten und darstellen \\
II & comparar & Gemeinsamkeiten/Unterschiede gewichtend gegenüberstellen \\
\rowcolor{lightgray}
II & explicar & Sachverhalte durch Wissen erläutern \\
\hline
III & comentar & Persönliche, begründete Stellungnahme abgeben \\
\rowcolor{lightgray}
III & discutir & Pro- und Kontra-Argumente abwägen und erörtern \\
III & evaluar & Nach Kriterien beurteilen und bewerten \\
\rowcolor{lightgray}
III & justificar & Aussagen durch Argumente schlüssig begründen \\
\hline
\end{tabularx}
\end{center}

\section{Stilmittel für Textanalyse}

\begin{center}
\begin{tabular}{l|l|l}
\rowcolor{lightblue}
\textbf{Deutsch} & \textbf{Spanisch} & \textbf{Ejemplo} \\
\hline
Alliteration & la aliteración & \textit{A las aladas almas...} \\
\rowcolor{lightgray}
Antithese & la antítesis & \textit{tan corto el amor, tan largo el olvido} \\
Metapher & la metáfora & \textit{Tus ojos son el mar} \\
\rowcolor{lightgray}
Personifikation & la personificación & \textit{Las estrellas nos miraban} \\
Vergleich & la comparación & \textit{hermoso como la noche} \\
\rowcolor{lightgray}
Wiederholung & la repetición & \textit{Temprano madrugó la madrugada} \\
Hyperbel & la hipérbole & \textit{Te lo he dicho mil veces} \\
\rowcolor{lightgray}
Ironie & la ironía & \textit{¡Qué bonito!} (negativ) \\
\hline
\end{tabular}
\end{center}

% ============================================================================
% TEIL V: EXAMENSHINWEISE
% ============================================================================
\newpage
\part*{Teil V: Examenshinweise mit Beispielantworten}
\addcontentsline{toc}{part}{Teil V: Examenshinweise}

\section{Prüfungsfragen mit Beispielantworten}

% === FRAGE 1 ===
\begin{exambox}[Frage 1: Rahmenplan]
\textbf{„Erläutern Sie die Unterschiede zwischen RP Gymnasium und RP Regionale Schule 2025."}
\end{exambox}

\textbf{Beispielantwort A (strukturiert):}
\begin{itemize}[leftmargin=0.5cm, itemsep=1pt]
\item \textbf{Stundenkontingent:} Gym hat in Jg. 7/8 je 120 USt, RegS nur 90 USt
\item \textbf{Zielniveau:} Gym erreicht Ende Jg. 10 Niveau B1, RegS nur A2+
\item \textbf{Progression:} Gym erreicht A2 bereits in Jg. 8, RegS erst in Jg. 9
\item \textbf{Anschlussfähigkeit:} Gym führt zur Qualifikationsphase, RegS zu MR/Berufsreife
\item \textbf{Inhalte:} Identische Themen, aber unterschiedliche Bearbeitungstiefe
\end{itemize}

\textbf{Beispielantwort B (kompakt):}
„Beide Rahmenpläne teilen dieselben verbindlichen Inhalte und die 80/20-Struktur. Der zentrale Unterschied liegt im Stundenkontingent -- das Gymnasium hat in Jg. 7 und 8 jeweils 30 Stunden mehr -- und im Zielniveau: B1 am Gymnasium versus A2+ an der Regionalen Schule. Dies erfordert eine angepasste Progression und Aufgabenkomplexität."

\vspace{0.4cm}

% === FRAGE 2 ===
\begin{exambox}[Frage 2: Kompetenzorientierung]
\textbf{„Wie fördern Sie die Sprachmittlungskompetenz in Jahrgangsstufe 9?"}
\end{exambox}

\textbf{Beispielantwort A (methodisch):}
\begin{itemize}[leftmargin=0.5cm, itemsep=1pt]
\item \textbf{Authentische Situationen:} Deutsche Touristen in Barcelona brauchen Hilfe beim Restaurantbesuch
\item \textbf{Textauswahl:} Deutsche Artikel zu Barcelona/Gaudí für spanischsprachige Austauschpartner zusammenfassen
\item \textbf{Scaffolding:} Redemittel-Karten, Wortfeldarbeit vorab, Modelltexte
\item \textbf{Progression:} Von mündlicher ad-hoc-Mittlung zu schriftlicher Zusammenfassung
\end{itemize}

\textbf{Beispielantwort B (mit konkretem Beispiel):}
„Im Themenbereich ‚Barcelona' lasse ich SuS einen deutschen Wikipedia-Artikel über Gaudí für eine fiktive spanische Schulklasse mitteln. Dabei geht es nicht um Übersetzung, sondern um adressatengerechte Auswahl und Neuformulierung. Vorbereitend erarbeiten wir Redemittel wie ‚El artículo trata de...', ‚Lo más importante es...' Die Aufgabe ist differenziert: Stärkere SuS mitteln frei, schwächere erhalten ein Strukturraster."

\vspace{0.4cm}

% === FRAGE 3 ===
\begin{exambox}[Frage 3: Differenzierung GK/LK]
\textbf{„Wie differenzieren Sie zwischen Grundkurs und Leistungskurs im Themenfeld 1 (El individuo y la sociedad)?"}
\end{exambox}

\textbf{Beispielantwort A (quantitativ + qualitativ):}
\begin{itemize}[leftmargin=0.5cm, itemsep=1pt]
\item \textbf{Zeitlich:} GK 45 USt vs. LK 75 USt -- mehr Vertiefung möglich
\item \textbf{Textumfang:} GK: Kurzfilm „Física II" + Interview; LK: zusätzlich Theaterstück „El método Grönholm"
\item \textbf{Anforderungen:} GK: Tagebucheintrag schreiben; LK: Rezension + Atmosphäre-Analyse
\item \textbf{AFB-Verteilung:} LK mit höherem AFB-III-Anteil (25--35\% statt 20--30\%)
\end{itemize}

\textbf{Beispielantwort B (am Beispiel):}
„Beim Kurzfilm ‚Física II' analysieren beide Kurse Jorges Beweggründe (AFB II). Der LK geht darüber hinaus: Er analysiert die filmischen Mittel zur Atmosphäregestaltung in der Bankszene und schreibt eine Rezension statt nur eines Tagebucheintrags. Zudem liest der LK ergänzend ‚El método Grönholm', was komplexere Bezüge zur Arbeitswelt ermöglicht."

\vspace{0.4cm}

% === FRAGE 4 ===
\begin{exambox}[Frage 4: Prüfungsformate]
\textbf{„Beschreiben Sie den Ablauf einer Paarprüfung im Abitur."}
\end{exambox}

\textbf{Beispielantwort (vollständig):}
„Die Paarprüfung dauert maximal 30 Minuten und gliedert sich in drei Teile:

\textbf{Teil 1 -- Aufwärmphase (ca. 2 Min.):} Kurzes Interview zum Ankommen. Geht \textit{nicht} in die Bewertung ein.

\textbf{Teil 2 -- Monologisches Sprechen (ca. 4 Min./Person):} Impulsgesteuertes Sprechen zu Bildern, Cartoons oder Grafiken. Jeder Prüfling spricht einzeln.

\textbf{Teil 3 -- Dialogisches Sprechen (ca. 5 Min./Person):} Situativ eingebettete Diskussion zu einem \textit{anderen} Themenschwerpunkt als Teil 2. Die Prüflinge interagieren miteinander.

\textbf{Bewertung (10 Min.):} 60\% Sprache (Aussprache, Wortschatz, Grammatik, Flüssigkeit), 40\% Inhalt (Aufgabenerfüllung, Argumentation, Interaktion). Zwei Lehrkräfte: Gesprächsleitung und Protokollant."

\vspace{0.4cm}

% === FRAGE 5 ===
\begin{exambox}[Frage 5: Operatoren]
\textbf{„Nennen Sie drei AFB-III-Operatoren mit Definitionen."}
\end{exambox}

\textbf{Beispielantwort:}
\begin{enumerate}[leftmargin=0.5cm, itemsep=2pt]
\item \textbf{comentar:} Eine persönliche, begründete Stellungnahme zu einem Sachverhalt, einer Problemstellung oder These abgeben. Erfordert eigene Wertung.
\item \textbf{discutir:} Eine Problemstellung erörtern, indem Pro- und Kontra-Argumente abgewogen und zu einem begründeten Urteil geführt werden.
\item \textbf{evaluar:} Sachverhalte, Aussagen oder Positionen nach explizit genannten oder selbst gewählten Kriterien beurteilen und bewerten.
\end{enumerate}

\textbf{Alternative Operatoren:} \textit{justificar} (Aussagen durch Argumente schlüssig begründen), \textit{valorar} (kritisch Stellung nehmen)

\vspace{0.4cm}

% === FRAGE 6 ===
\begin{exambox}[Frage 6: Literaturauswahl]
\textbf{„Welche Lektüre würden Sie für das Thema Migration wählen und warum?"}
\end{exambox}

\textbf{Beispielantwort A (für GK):}
„Für den Grundkurs wähle ich \textbf{‚Abdel'} von Enrique Páez. Der Roman ist sprachlich zugänglich (ca. A2/B1), erzählt chronologisch und bietet starke Identifikationsmöglichkeiten: Der 16-jährige Protagonist entspricht dem Alter der SuS. Die Themen illegale Einwanderung, Fluchtursachen und Integration in Spanien sind hochaktuell und ermöglichen Bezüge zur europäischen Flüchtlingsdebatte."

\textbf{Beispielantwort B (für LK):}
„Im Leistungskurs setze ich \textbf{‚La casa en Mango Street'} von Sandra Cisneros ein. Die Vignettenstruktur ermöglicht flexible Kapitelauswahl und literarische Analyse (Metaphorik, Erzählperspektive). Der Text behandelt Chicano-Identität, was Varietäten des Español Americano thematisiert. Zusätzlich nutze ich Interviews mit der Autorin für AFB-III-Aufgaben zur Autorenintention."

\textbf{Beispielantwort C (vergleichend):}
„Die Wahl hängt vom Kursniveau ab: ‚Abdel' für GK wegen Zugänglichkeit, ‚La casa en Mango Street' für LK wegen literarischer Komplexität. Beide Texte thematisieren Migration, aber aus unterschiedlichen Perspektiven: Abdel als Betroffener in Spanien, Esperanza als Chicana in den USA."

\vspace{0.4cm}

% === FRAGE 7 ===
\begin{exambox}[Frage 7: Sprachbildung]
\textbf{„Wie integrieren Sie sprachbildende Maßnahmen in Ihren Spanischunterricht?"}
\end{exambox}

\textbf{Beispielantwort A (nach RP Sprachbildung):}
\begin{itemize}[leftmargin=0.5cm, itemsep=1pt]
\item \textbf{Wortschatzarbeit:} Wortfelder visualisieren, Kollokationen üben, Wortfamilien erschließen
\item \textbf{Textmuster:} Explizite Vermittlung von Textsortenmerkmalen (Brief, Kommentar, Analyse)
\item \textbf{Scaffolding:} Satzanfänge, Redemittel-Karten, Strukturierungshilfen
\item \textbf{Sprachreflexion:} Vergleich Spanisch--Deutsch, Interferenzen thematisieren
\item \textbf{Operatorentraining:} Bedeutung und sprachliche Umsetzung von Operatoren üben
\end{itemize}

\textbf{Beispielantwort B (konkret):}
„Vor einer Schreibaufgabe ‚comentario' erarbeite ich mit den SuS explizit den Textaufbau (Einleitung--Hauptteil--Schluss), typische Redemittel (‚En mi opinión...', ‚Por un lado... por otro lado...') und übe die Bedeutung des Operators ‚comentar'. Dies ist Sprachbildung im Fach: Die SuS lernen nicht nur Spanisch, sondern auch, wie man strukturiert argumentiert."

\vspace{0.4cm}

% === FRAGE 8 ===
\begin{exambox}[Frage 8: Digitalität und KI]
\textbf{„Wie setzen Sie KI-Tools reflektiert im Spanischunterricht ein?"}
\end{exambox}

\textbf{Beispielantwort A (kritisch-konstruktiv):}
„Der RP 2025 betont, dass KI-Tools \textit{nicht ohne Reflexion} eingesetzt werden sollen. Mein Ansatz:
\begin{itemize}[leftmargin=0.5cm, itemsep=1pt]
\item \textbf{Transparenz:} Mit SuS besprechen, was LLMs können und was nicht (Halluzinationen, keine Quellenprüfung)
\item \textbf{Nicht im Pre-Writing:} Brainstorming und Strukturierung selbst leisten
\item \textbf{Sinnvoller Einsatz:} Überarbeitungsphase, Grammatikcheck, Ideenerweiterung nach eigener Vorarbeit
\item \textbf{Kritische Prüfung:} KI-generierte Texte gemeinsam auf Fehler und Floskeln analysieren
\end{itemize}"

\textbf{Beispielantwort B (praktisches Beispiel):}
„Ich lasse SuS einen eigenen Kommentar schreiben und dann von einer KI ‚verbessern'. Anschließend vergleichen wir kritisch: Was hat die KI geändert? Ist es wirklich besser? Hat sie den Inhalt verändert? So entwickeln SuS Urteilsfähigkeit und erkennen, dass KI ein Werkzeug ist, das menschliche Reflexion nicht ersetzt."

\textbf{Beispielantwort C (ablehnend-begründet):}
„In der Textproduktion setze ich KI bewusst nicht ein, da der Schreibprozess selbst lernförderlich ist. Beim eigenständigen Formulieren werden Wortschatz aktiviert und grammatische Strukturen gefestigt. KI-generierte Texte entziehen den SuS diese Lernchance. Allenfalls zur Fehleranalyse im Nachgang ist ein reflektierter Einsatz denkbar."

\vspace{0.8cm}

\section{Fallbeispiel: Unterrichtsplanung}

\begin{fallbeispielbox}[Fallbeispiel: Einheit „Movimientos migratorios" (Thema 3, LK)]
\textbf{Kontext:} LK Spanisch, Jg. 12, 14 SuS, ca. 20 USt

\textbf{Sequenzierung:}
\begin{enumerate}[leftmargin=0.5cm, itemsep=2pt]
\item \textbf{Einstieg (2 USt):} Doku „México--EE.UU.: El gran cruce" -- Hörsehverstehen
\item \textbf{Erarbeitung I (4 USt):} „La casa en Mango Street" -- Kapitelarbeit, Leseverstehen
\item \textbf{Vertiefung (3 USt):} Interview Sandra Cisneros -- Autorenintention, Analyse
\item \textbf{Produktion (3 USt):} Brief aus Esperanzas Sicht -- Schreiben
\item \textbf{Transfer (4 USt):} Aktuelle Migration in Europa -- Sprachmittlung, Diskussion
\item \textbf{Erarbeitung II (2 USt):} Varietäten des Español Americano -- Sprachbewusstheit
\item \textbf{Abschluss (2 USt):} Zusätzliches Kapitel verfassen -- kreatives Schreiben
\end{enumerate}

\textbf{Leistungsüberprüfung:} Klausur (Textanalyse + comentario) ODER Sprechprüfung (Diskussion)

\textbf{Differenzierung LK:}
\begin{itemize}[leftmargin=0.5cm, itemsep=1pt]
\item Zusätzliche Texte (Interview, weitere Kapitel)
\item Höhere Anforderungen bei Analyse (Gestaltungsmittel)
\item Komplexere Schreibaufgaben (eigenes Kapitel)
\end{itemize}
\end{fallbeispielbox}

% ============================================================================
% TEIL VI: CHECKLISTEN
% ============================================================================
\newpage
\part*{Teil VI: Checklisten}
\addcontentsline{toc}{part}{Teil VI: Checklisten}

\begin{checklistbox}[Checkliste: Vor dem Unterricht (Sek I)]
\begin{itemize}[leftmargin=0.5cm, label=\checkbox]
\item Richtigen Rahmenplan identifiziert? (RP 2002 vs. RP 2025)
\item Stundenkontingent beachtet? (Gym 120/90 USt vs. RegS 90 USt)
\item GER-Zielniveau bekannt? (Gym: A1$\rightarrow$B1 / RegS: A1$\rightarrow$A2+)
\item Verbindliche Inhalte der Jahrgangsstufe geprüft?
\item RP Sprachbildung als Ergänzung berücksichtigt?
\item 80/20-Regel: Freiraum für eigene Schwerpunkte eingeplant?
\end{itemize}
\end{checklistbox}

\vspace{0.5cm}

\begin{checklistbox}[Checkliste: Vor dem Unterricht (Sek II)]
\begin{itemize}[leftmargin=0.5cm, label=\checkbox]
\item Themenfeld identifiziert (1--4)?
\item GK/LK-Unterschiede beachtet (USt, Niveau, Zusatzanforderungen)?
\item Operatorenliste griffbereit?
\item Literatur ausgewählt und beschafft?
\item Verknüpfungsbeispiele des IQ M-V genutzt?
\item Mind. 1 komplexe Leistungsermittlung Sprechen in QPhase eingeplant?
\item Vorabhinweise für aktuelles Abiturjahr geprüft?
\end{itemize}
\end{checklistbox}

\vspace{0.5cm}

\begin{checklistbox}[Checkliste: Vor der Sprechprüfung]
\begin{itemize}[leftmargin=0.5cm, label=\checkbox]
\item Zeitplan erstellt (30 Min./Tandem oder 40 Min./3er-Gruppe)?
\item Räume reserviert (Prüfungsraum + Aufenthaltsraum)?
\item Jury organisiert (2 Fachlehrkräfte)?
\item Impulsmaterial vorbereitet (Bilder, Cartoons, Grafiken)?
\item Teil 2 und Teil 3 mit VERSCHIEDENEN Themenschwerpunkten?
\item Bewertungsbögen ausgedruckt?
\item Regieanweisungen (Interlocutor Frame) vorbereitet?
\item Protokollformular bereit?
\end{itemize}
\end{checklistbox}

\vspace{0.5cm}

\begin{checklistbox}[Checkliste: Klausurerstellung (Abiturniveau)]
\begin{itemize}[leftmargin=0.5cm, label=\checkbox]
\item AFB-Verteilung korrekt? (20\% / 45\% / 35\%)
\item Operatoren aus verbindlicher Liste verwendet?
\item Textlänge angemessen? (GK: 400--600 Wörter / LK: 600--800)
\item Aufgabenstellung eindeutig formuliert?
\item Erwartungshorizont erstellt?
\item Bewertungsraster vorhanden (Sprache/Inhalt)?
\item Stilmittel-Analyse bei literarischen Texten gefordert?
\end{itemize}
\end{checklistbox}

% ============================================================================
% RESSOURCEN-ANHANG
% ============================================================================
\newpage
\part*{Anhang: Ressourcen}
\addcontentsline{toc}{part}{Anhang: Ressourcen}

\section{Vollständige Literaturliste nach Themen}

\subsection{Thema 1: El individuo y la sociedad}
\begin{itemize}[leftmargin=0.5cm, itemsep=2pt]
\item \textbf{Desconexión total} (Lourdes Miquel, 2020) -- Roman, Gruppenarbeit
\item \textbf{La Casa} (Paco Roca, 2015) -- Novela Gráfica, Astiberri
\item \textbf{La mirada} (Juan Madrid, 1991) -- Kurzgeschichte, in: Cuentos del asfalto
\item \lmark{} \textbf{El método Grönholm} (Jordi Galceran, 2003) -- Theaterstück, kostenlos annotiert
\end{itemize}

\subsection{Thema 2: La identidad nacional}
\begin{itemize}[leftmargin=0.5cm, itemsep=2pt]
\item \textbf{El desertador} (José María Merino, 1941) -- Kurzgeschichte, Bürgerkrieg
\item \textbf{Una fosa poco común} (Luis García Jambrina, 2009) -- Krimi, Massengrab
\item \textbf{Las venas abiertas de América Latina} (Eduardo Galeano, 2003) -- Sachbuch
\end{itemize}

\subsection{Thema 3: Aspectos de política}
\begin{itemize}[leftmargin=0.5cm, itemsep=2pt]
\item \textbf{La piel de un indio no cuesta caro} (Julio Ramón Ribeyro, 1961) -- Kurzgeschichte
\item \textbf{Abdel} (Enrique Páez, 1994) -- Roman, Migration
\item \textbf{La casa en Mango Street} (Sandra Cisneros) -- Roman, Chicano-Literatur
\end{itemize}

\subsection{Thema 4: Desafíos modernos}
\begin{itemize}[leftmargin=0.5cm, itemsep=2pt]
\item \textbf{547 amigos} (Lorenzo Silva) -- Kurzgeschichte, soziale Netzwerke
\item \textbf{El agua y la humanidad} (Rigoberta Menchú) -- Sachtext, Umwelt
\item \lmark{} \textbf{El guardagujas} (Juan José Arreola, 1952) -- Kurzgeschichte, fantastisch
\end{itemize}

\section{URL-Übersicht}

\begin{center}
\footnotesize
\begin{tabularx}{\textwidth}{l|X}
\rowcolor{lightblue}
\textbf{Ressource} & \textbf{URL} \\
\hline
Rahmenpläne MV & \url{https://www.bildung-mv.de/unterricht/rahmenplaene/} \\
\rowcolor{lightgray}
El método Grönholm & \url{https://www.bildung-mv.de/.../spanisch/} \\
Kurzfilm Física II & \url{https://cortosdemetraje.com/fisica-ii/} \\
\rowcolor{lightgray}
Castro-Rede Video & \url{https://www.youtube.com/watch?v=8eX8cbT1Tfo} \\
Castro-Rede Text & \url{https://desinformemonos.org/discurso-fidel-castro-frente-la-onu-1979/} \\
\rowcolor{lightgray}
Menchú Video & \url{https://www.youtube.com/watch?v=j4EUXUmJML8} \\
Menchú Text & \url{https://www.zaragoza.es/.../Menchu_ES.pdf} \\
\rowcolor{lightgray}
547 amigos & \url{https://www.lorenzo-silva.com/archivo/547amigos.html} \\
El desertador & \url{https://narrativabreve.com/.../desertor.html} \\
\rowcolor{lightgray}
Galeano Volltext & \url{https://static.telesurtv.net/.../las_venas_abiertas.pdf} \\
Ribeyro & \url{https://www.literatura.us/julio/indio.html} \\
\hline
\end{tabularx}
\end{center}

% ============================================================================
% QUICK-REFERENCE (LETZTE SEITE)
% ============================================================================
\newpage
\section*{QUICK-REFERENCE: Spanisch MV auf einen Blick}
\addcontentsline{toc}{section}{Quick-Reference (1 Seite)}

\begin{center}
\fbox{\parbox{0.95\textwidth}{
\centering\Large\bfseries\color{mainblue} SPANISCH MECKLENBURG-VORPOMMERN -- QUICK-REFERENCE
}}
\end{center}

\vspace{0.3cm}

\begin{minipage}[t]{0.48\textwidth}
\textbf{\color{mainblue}RAHMENPLÄNE}
\begin{itemize}[leftmargin=0.3cm, itemsep=1pt, topsep=2pt]
\item \textbf{Sek I Gym:} RP 2025 (ab 25/26, aufsteigend)
\item \textbf{Sek I RegS:} RP 2025 (ab 25/26, aufsteigend)
\item \textbf{Sek II:} RP 2019 (unverändert gültig)
\item \textbf{RP 2002:} auslaufend bis 2028/29
\end{itemize}

\vspace{0.2cm}
\textbf{\color{mainblue}NIVEAUSTUFEN (Ende Jg. 10)}
\begin{itemize}[leftmargin=0.3cm, itemsep=1pt, topsep=2pt]
\item Gymnasium: \textbf{B1}
\item Regionale Schule: \textbf{A2+}
\item Sek II (fortgef.): \textbf{B2}
\end{itemize}

\vspace{0.2cm}
\textbf{\color{mainblue}4 THEMENFELDER SEK II}
\begin{enumerate}[leftmargin=0.5cm, itemsep=1pt, topsep=2pt]
\item El individuo y la sociedad
\item La identidad nacional
\item Aspectos de política
\item Desafíos modernos
\end{enumerate}

\vspace{0.2cm}
\textbf{\color{mainblue}PFLICHTPRÜFUNG}
\begin{itemize}[leftmargin=0.3cm, itemsep=1pt, topsep=2pt]
\item Mind. 1× Sprechen in QPhase
\item Paarprüfung ab 2025 möglich
\item 60\% Sprache / 40\% Inhalt
\end{itemize}
\end{minipage}
\hfill
\begin{minipage}[t]{0.48\textwidth}
\textbf{\color{mainblue}AFB-VERTEILUNG ABITUR}
\begin{itemize}[leftmargin=0.3cm, itemsep=1pt, topsep=2pt]
\item AFB I: 20\% (describir, resumir, presentar)
\item AFB II: 45\% (analizar, comparar, explicar)
\item AFB III: 35\% (comentar, discutir, evaluar)
\end{itemize}

\vspace{0.2cm}
\textbf{\color{mainblue}PAARPRÜFUNG (30 MIN.)}
\begin{enumerate}[leftmargin=0.5cm, itemsep=1pt, topsep=2pt]
\item Aufwärmphase: 2 Min. (keine Wertung)
\item Monolog: 4 Min./Person
\item Dialog: 5 Min./Person
\item Bewertung: 10 Min.
\end{enumerate}

\vspace{0.2cm}
\textbf{\color{mainblue}WICHTIGE DOKUMENTE}
\begin{itemize}[leftmargin=0.3cm, itemsep=1pt, topsep=2pt]
\item Operatorenliste (verbindlich!)
\item HR Sprechprüfung
\item HR Paarprüfung (ab 2025)
\item Verknüpfungsbeispiele (4×)
\item Literaturempfehlungen
\end{itemize}

\vspace{0.2cm}
\textbf{\color{mainblue}KONTAKT}\\
IQ M-V, Fachbereich 4\\
Tel: 0385 588 17003\\
\url{bildung-mv.de/rahmenplaene}
\end{minipage}

\vspace{0.5cm}
\begin{center}
\begin{tikzpicture}
\node[rectangle, rounded corners=3pt, fill=lightgray, draw=gray, minimum width=14cm, minimum height=0.8cm] {
  \small\textit{Erstellt auf Basis von 22 offiziellen Dokumenten des IQ M-V | Stand: Dezember 2025 | v4}
};
\end{tikzpicture}
\end{center}

\end{document}
